\section{Conclusion}

The retained Blackwell backward uses adaptive FP4 Q/K for score and
$dQ/dK$, E4M3 operands for $dP/dV$, direct E4M3 dS production, FP32 TMEM
accumulation, and BF16 dQ reduction.  Its represented equations are explicit:
P is scaled by $2^8$, $dP-r$ by $2^4$, and their product lands directly at
the $2^{12}$ dS scale consumed by mixed E4M3$\times$E2M1 tensor operations.
This route is 11.3--16.4\% faster than the copied BF16 control on the measured
long-sequence matrix while preserving approximately 0.9915 dQ/dK cosine in
the calibrated regime.

The native NVFP4 learned QKV projection is the larger systems result.  It
emits one forward/backward Q/K payload, exact backward views, transposed
MXFP4 V, and optional BF16 Q/K/V from its output fragments.  On the measured
projection shapes it is 1.41--1.68$\times$ faster than three BF16 GEMMs plus
Q/K packing with BF16 publication, and 1.60--1.93$\times$ faster without it.
Pair-native RoPE is fused before every Q/K publication and saves
16.128~$\mu$s against a strong standalone forward transform at
S4096/H24/K3072.  On the final five-shape matrix that keeps FP4 FA4 forward in
every low-precision row, the RoPE-aware dense decoder boundary is
1.118--1.720$\times$ faster than BF16 for activation gradients and
1.096--1.469$\times$ when projection weight gradients are charged.  At
S4096/H24/K3072, full-training algorithmic MFU rises from 26.94\% to 32.38\%
on a common 2.5-PFLOP/s dense denominator.  This boundary remains a dense
$D_{QK}=192$ result, not yet a stock grouped-query Llama drop-in.  The fused
dO row/column/delta producer adds a separate 1.935$\times$ setup win.

The final stacked QKV-gradient producer fuses inverse RoPE, K16 scale
selection, and NVFP4 packing before one dX projection.  With the normal
unsignaled dQ endpoint it improves the S4096/H24 activation boundary from
1.578336 to 1.566048~ms; a genuinely hierarchical two-lane producer regresses
to 1.635936~ms because its readiness signaling slows attention.  The stacked
winner is therefore activation-only.  Full training keeps the interleaved
BF16 QKV gradient needed for weight gradients.

A 500-step, three-seed teacher/student proxy provides the first optimization
trajectory rather than another one-step cosine.  Low precision retains 93.2\%
of BF16's absolute validation-loss improvement on average and reduces its
floor-adjusted excess error at least as quickly in all three runs.  Raw native
loss is still higher because the FP4 forward has a 0.093--0.098 relative-MSE
teacher floor, and a common BF16 evaluator leaves final loss 25.5\% higher on
average.  These data support similar block-level optimization progress, not
equal final loss or stock-Llama language-model convergence.

For the aggressive pure-FP4 endpoint, the projection epilogue now performs a
single fixed-scale Q/K quantization and publishes all six required backward
layouts directly.  This reduces unified projection time by 6.8--14.6\% and
the measured projection-plus-pure-backward chains by 6.7--9.5\%.  The pure
kernel is 10.0--11.4\% faster than BF16 on the seven-shape sweep, but remains
slower than the mixed-dP hybrid and its dQ/dK cosine is only 0.864--0.887.
It is therefore a completed aggressive endpoint, not the new default.
Substituting NVFP4 K16 scales does not rescue it: NVFP4 dS improves component
cosine to approximately 0.969 only by becoming slower than BF16, NVFP4 Q/K
regresses both speed and accuracy, and the combined diagnostic is non-finite.
Pure FP4 is therefore not a competitive backward default in the current
SM100 ownership and scale-publication topology.

The repaired adaptive mixed-dP route is a valid intermediate Pareto point.
It is 2.3--2.7\% faster than adaptive FP8 dP/dV at S4096/H24, but 1.3--1.7\%
slower at S8192/H8 and lowers calibrated dQ/dK cosine from approximately
0.991 to 0.983.  It is exposed explicitly rather than selected automatically.

The integration experiments also establish the present ceiling.  Local
represented-P reuse is worth only 0.76\%; an HBM P cache is too large; and the
96-byte mixed-dP transport is not a legal completing transaction.  Moving
per-head dQ readiness into the projection K loop, ordering the async proxy,
and alternating two private TMEM accumulators improves the original direct
consumer enough to beat attention plus cuBLAS by 2.6~$\mu$s at S8192/H8.
The same bounded topology remains 7.9~$\mu$s behind at S4096/H24 even after a
vendor-GEMM tail, so wider heads shape-dispatch to the non-regressing BF16
dQ plus cuBLAS chain.  The next breakthrough is therefore not another barrier
rearrangement.  It is projection-native adaptive scale production and a
reduction topology that lets projection backward consume dQ before a
standalone global BF16 tensor exists.  A direct two-lane implementation now
proves that topology algebraically and numerically, but remains 32.1~$\mu$s
behind at S8192/H8 because it preserves BF16 partial traffic and duplicates
the projection product.  The remaining requirement is stricter: group owners
on chip before publication, or publish compact FP8/FP4 partials whose repeated
products are paid for by low-precision tensor throughput.

The compact NVFP4 final-tile experiment validates the consumer half of that
direction.  With delayed device-resident scale state, NVFP4 conversion plus
projection reaches approximately 0.991 cosine and the composed H24 chain is
3.7\% faster than BF16; H8 is 1.9\% slower.  It still quantizes completed
private BF16 dQ scratch.  The next implementation must move that quantizer to
tile readiness or an on-chip super-owner so the low-precision handoff removes,
rather than follows, the BF16 reduction round trip.  The analogous RoPE
opportunity has now been tested.  Fusing $R^{\mathsf T}dQ$ and
$R^{\mathsf T}dK$ into the delayed-scale NVFP4 quantizer is bit-exact, but
saves only 2.720~$\mu$s without BF16 republication and 3.392~$\mu$s with it.
At S4096/H24/K4096 those component savings become 0.46\% for activation
gradients and 0.60\% when both projection weight-gradient GEMMs are charged.
The row-owned quantizer cannot make the 0.115-ms transform free.  A
forward-sized systems gain requires inverse RoPE inside a projection consumer
that receives the completed reduction tile before global BF16 publication.
